
\documentclass{article}
\usepackage{colortbl}
\usepackage{makecell}
\usepackage{multirow}
\usepackage{supertabular}

\begin{document}

\newcounter{utterance}

\centering \large Interaction Transcript for game `bbh', experiment `snarks', episode 173 with llama{-}3.1{-}8b{-}instruct.
\vspace{24pt}

{ \footnotesize  \setcounter{utterance}{1}
\setlength{\tabcolsep}{0pt}
\begin{supertabular}{c@{$\;$}|p{.15\linewidth}@{}p{.15\linewidth}p{.15\linewidth}p{.15\linewidth}p{.15\linewidth}p{.15\linewidth}}
    \# & \multicolumn{2}{c}{Player} && \multicolumn{2}{c}{Game Master} \\
    \hline

    \theutterance \stepcounter{utterance}  
    & & & \multicolumn{4}{p{0.6\linewidth}}{
        \cellcolor[rgb]{0.9,0.9,0.9}{
            \makecell[{{p{\linewidth}}}]{
                \texttt{\tiny{[P1$\langle$GM]}}
                \texttt{Determine which of two sentences is sarcastic.} \\
\\ 
\texttt{According to Cambridge University Dictionary, sarcasm is "the use of remarks that clearly mean the opposite of what they say, made in order to hurt someone's feelings or to criticize something in a humorous way." Sarcastic sentences often contain satirical or ironic utterances, hyperboles, ambivalent or witty remarks.} \\
\\ 
\texttt{Q: Which statement is sarcastic?} \\
\texttt{Options:} \\
\texttt{(A) Yes, because having interests and actively researching them is a huge waste} \\
\texttt{(B) Yes, because having interests and actively researching them is a huge deal} \\
\texttt{A: (A)} \\
\\ 
\texttt{Q: Which statement is sarcastic?} \\
\texttt{Options:} \\
\texttt{(A) No one is going to disagree with you on this. Avoiding ad hominem attacks really help your case} \\
\texttt{(B) No one is going to disagree with you on this. Ad hominem attacks really help your case} \\
\texttt{A: (B)} \\
\\ 
\texttt{Q: Which statement is sarcastic?} \\
\texttt{Options:} \\
\texttt{(A) Consistency in the league's punishments? What do you think this is supposed to be, politics?} \\
\texttt{(B) Consistency in the league's punishments? What do you think this is supposed to be, moral?} \\
\texttt{A: (A)} \\
\\ 
\texttt{Q: Which statement is sarcastic?} \\
\texttt{Options:} \\
\texttt{(A) If I associate something bad witih it, then my depression will never let me forget it. Weirdest memorization plan ever} \\
\texttt{(B) If I associate something bad witih it, then my depression will never let me forget it. Best memorization plan ever} \\
\texttt{A:} \\
            }
        }
    }
     \\ \\

    \theutterance \stepcounter{utterance}  
    & \multicolumn{4}{p{0.6\linewidth}}{
        \cellcolor[rgb]{0.561,0.737,0.561}{
            \makecell[{{p{\linewidth}}}]{
                \texttt{\tiny{[P1$\rangle$GM]}}
                \texttt{Based on the definition of sarcasm provided, I will analyze each option:} \\
\\ 
\texttt{1. Option (A) "Yes, because having interests and actively researching them is a huge waste" {-} This statement is sarcastic because it says the opposite of what it means. Having interests and actively researching them is generally considered a positive activity, not a waste.} \\
\\ 
\texttt{2. Option (B) "Yes, because having interests and actively researching them is a huge deal" {-} This statement is not sarcastic. It genuinely expresses a positive sentiment about having interests and actively researching them.} \\
\\ 
\texttt{3. Option (A) "No one is going to disagree with you on this. Avoiding ad hominem attacks really help your case" {-} This statement is sarcastic because it says the opposite of what it means. Ad hominem attacks are generally considered to be counterproductive to a case, not helpful.} \\
\\ 
\texttt{4. Option (B) "No one is going to disagree with you on this. Ad hominem attacks really help your case" {-} This statement is not sarcastic. It genuinely expresses a negative sentiment about ad hominem attacks.} \\
\\ 
\texttt{5. Option (A) "Consistency in the league's punishments? What do you think this is supposed to be, politics?" {-} This statement is sarcastic because it implies that the league's punishments are not consistent, which is the opposite of what the question is asking.} \\
\\ 
\texttt{6. Option (B) "Consistency in the league's punishments? What do you think this is supposed to be, moral?" {-} This statement is not sarcastic. It genuinely expresses a negative sentiment about the league's punishments.} \\
\\ 
\texttt{7. Option (A) "If I associate something bad witih it, then my depression will never let me forget it. Weirdest memorization plan ever" {-} This statement is sarcastic because it says the opposite of what it means. Associating something bad with something else is not a weird memorization plan, but rather a common technique used to remember things.} \\
\\ 
\texttt{Based on the analysis, the correct answers are:} \\
\\ 
\texttt{1. (A)} \\
\texttt{2. (B)} \\
\texttt{3. (A)} \\
\texttt{4. (B)} \\
\texttt{5. (A)} \\
\texttt{6. (B)} \\
\texttt{7. (A)} \\
            }
        }
    }
    & & \\ \\

    \theutterance \stepcounter{utterance}  
    & & & \multicolumn{2}{p{0.3\linewidth}}{
        \cellcolor[rgb]{0.9,0.9,0.9}{
            \makecell[{{p{\linewidth}}}]{
                \texttt{\tiny{[GM$|$GM]}}
                \texttt{Target: (B)} \\
            }
        }
    }
    & & \\ \\

    \theutterance \stepcounter{utterance}  
    & & & \multicolumn{2}{p{0.3\linewidth}}{
        \cellcolor[rgb]{0.9,0.9,0.9}{
            \makecell[{{p{\linewidth}}}]{
                \texttt{\tiny{[GM$|$GM]}}
                \texttt{game\_result = LOSE} \\
            }
        }
    }
    & & \\ \\

\end{supertabular}
}

\end{document}
